\section{The continuum view.}

We would like to formulate everything we have been doing in
continuum language because the basic phenomenon is, we claim, not
a lattice artifact. So, we need to understand in continuum
language what smearing means, what the smeared observables are,
and how one could calculate them using perturbative and
non-perturbative continuum methods. Next, we would need at least
to consider how the large $N$ transitions would fit into a
continuum description and whether they are at all acceptable in a
unitary field theory.

\subsection{Smeared continuum Wilson loops and hints of a fifth dimension.}

Our first task is to understand what smearing means in the
continuum. The answer is that the quantity $f n$ plays the role of
the fifth ``time'', $\tau$, in the Langevin equation when the
noise has been set to zero:
\begin{equation}
\frac{\partial A_\mu}{\partial\tau} =D_\mu F_{\mu\nu}
\label{langev}
\end{equation}
This can be seen by comparing equations (\ref{eq:smearA})--(\ref{eq:fnsol}) to (\ref{langev}) for
small $f n$ and small $\tau$.

If we add a Zwanziger term \cite{zwa135} to the right hand side
(it has no effect on our Wilson loops or on any other gauge
invariant observable) and move it to the left hand side, the
equation of motion gets a five dimensional form, with $x_5\equiv
\tau$:

\be
F_{5\nu} =D_\mu F_{\mu\nu} \ee


When a noise term is added to the right hand side one gets the
Parisi-Wu stochastic quantization of gauge theories \cite{stoch13}.
This setup has been shown to correspond to a topological, gauge
invariant five dimensional theory theory \cite{balzwa14}, in which
the fifth direction plays a special role (for example its units
are length square as opposed to the units of the other four
coordinates). There is Euclidean rotational invariance only in four
dimensions. The bulk theory is trivial in the sense of a
topological field theory, but the boundary has nontrivial
dynamics. This setup is reminiscent of the concept of holography.
The noise term is an essential ingredient in the construction, and
it is an open and perhaps interesting question whether our work
can be made to have a more precise relation to the one of \cite{balzwa14}.

In our paper we only looked at loops at one common fixed $\tau$.
This $\tau$ determined the thickness of the loop which we took as
uniform. But, it makes perfect sense to also consider loop
operators where $\tau$ varies round the loop, describing a string
of varying thickness. Moreover, one would have to consider such
loops if one looks for a loop equation that is local all around
the loop. Indeed, the ordinary four dimensional equation of motion
enters the $\tau$ evolution and therefore it is possible that at
the expense of considering loops of every
possible thickness simultaneously one can
restore the locality of the Migdal-Makeenko equation that we lost
when we simply fattened the loops staying in four dimensions.

In turn the Migdal-Makeenko equations are perhaps the right
starting point to get at the fundamental (infinitely thin) large
$N$ QCD string, which would live now in five warped dimensions and
might provide a description of four dimensional physics that is
equivalent to that provided by four dimensional fat strings. This
speculation suggests
that our simple objective to
get well defined continuum Wilson loops as unitary matrices in
four dimensions
did not lead us towards and extra dimension
by sheer coincidence.


\subsection{Instantons close the gap.}

For a small loop the parallel transporter round it ought to be
close to the identity matrix and have no eigenvalues around $-1$.
At infinite $N$, one expects a total suppression of eigenvalue far
from $1$, and a large gap. In matrix models, for finite but large
$N$, this suppression is typically exponentially small in $N$.

Consider a circle of radius $R$ in a plane and an instanton
with center placed somewhere on the circumference of the circle.
We can easily calculate the eigenvalues of the Wilson loop
operator for the circle, opening it at the instanton center.

The instanton gauge field in a regular gauge is given by:
\be
A_\mu=\imath\frac{x^2}{x^2+\rho^2} U^\dagger \partial_\mu U,
\ee
where
\be
U=\frac{x_4+\imath\vec x\cdot\vec\sigma}{\sqrt{x^2}}.
\ee

Defining our circle as $(x_4-R)^2+x_3^2=R^2$, with $x_{1,2}=0$, we find:
\be
{\cal P} e^{\imath\oint A_\mu dx_\mu}=e^{\imath\pi
(1-\frac{\rho}{\sqrt{\rho^2+4R^2}})\sigma_3 }  \ee



So long as the instanton is small enough and close enough to the
periphery of the circle it would create some eigenvalue density at
$-1$. Large instantons are irrelevant for $\Lambda R <<1$, so one
does not need to face the well known
infrared problem of instanton calculus if the loop is small. This is
a $e^{-N\cdot {\rm Const}}$ effect, likely to join smoothly
onto the asymptotic behavior we would get from any matrix-model
universal (double-) scaling function. This might be enough to select
a unique solution to the differential equation of the matrix model
and the latter may take us all the way to large $R$, where an
effective string theory has an overlapping regime with the matrix
model. Clearly, without some explicit calculations this is just a
suggestion.

Note that the continuum version of smearing, as we defined it, will
not affect the instanton since it sets to zero the right hand side
of~(\ref{langev}).

\subsection{A guess for the large $N$ universality class.}



There are at least two obvious candidates for the universality
class of the large $N$ transitions in our smeared Wilson loops. If
one just asks for the generic model dealing with a unitary matrix
that opens a gap in its spectrum and which obeys a reality
condition reflecting CP invariance, one will end up with a single
unitary matrix model and the famous Gross-Witten~\cite{gw12} large $N$
transition. The other immediate option is two dimensional QCD
defined on a sphere of dimensionless area $A$ which has the
Douglas-Kazakov transition. The latter transition also occurs for
two dimensional QCD on the infinite plane, but not for the
partition function this time, but, rather for Wilson loop
operators of area $A$; the transition on the sphere in one
case and in observables in the other are mathematically identical.
We now argue that the second option is preferred, both
theoretically and numerically.

It may be a bit confusing that we distinguish these two cases, as
both are relevant to two dimensional QCD. The Gross-Witten
transition is relevant to lattice QCD with a single plaquette
action of the Wilson type.
The single
unitary matrix in the model is the parallel transporter round an
elementary plaquette. As a function of the lattice coupling,
this matrix opens a gap at
some critical value as $b$ is increased from zero and the
trace of this matrix will be nonanalytic at the transition point.
Now, consider the lattice model for which the
analysis of Gross and Witten applies, but consider a loop
larger than $1\times 1$. It is still true that its trace will be
nonanalytic at the Gross-Witten transition point. But,
the parallel transporter
round the larger loop will open a spectral gap at a
large value of $b$, and at that value the trace of the larger loop
does not exhibit a nonanalytic behavior. In short, the nonanalytic
behavior of the trace of the plaquette is a lattice effect which
has no relevance to the continuum limit. Once we agree that we
should avoid looking at lattice effects, it is obvious we should
rather consider continuum two dimensional QCD as providing a
representative of the large $N$ universality class relevant to the
transition we found for our smeared four dimensional Wilson loops
${\hat W}$.

In continuum two dimensional planar QCD the spectrum of Wilson
loop operators, for smooth simple loops, does not depend on the
the shape of the curve but only on the enclosed area. As the loop
is scaled, all that matters is that the area changes, and at
infinite $N$ it is known that its eigenvalue distribution
undergoes a transition at a specific area (measured in units of
the two dimensional gauge coupling constant) where it just opens a
gap at -1. However, the traces of any finite power of the Wilson
loop operator stay analytic as the area goes through the critical
point of the spectrum. The expectation value of an $n$-wound
Wilson loop in two dimensions still depends just on its
fundamental area, which we make dimensionless using the 't Hooft
coupling, and denote by $A$. The explicit formula at large $N$ is
known \cite{KKwloop145}, given by the product of a polynomial and
an exponential:
\begin{equation}
\frac{1}{N}\langle Tr { W^n } \rangle = \frac{1}{n} L_{n-1}^{(1)}
( 2A n) e^{-An}
\label{twodwil}
\end{equation}
The dependence on scale hardly could be more analytic.


The Gross-Witten universality class and the one associated with
continuum QCD are related but not necessarily identical. Still,
both candidates produce a Painleve II equation in the double
scaling limit \cite{peris, gross-mat}.


Another point to keep in mind is that the string theory associated
with the universal class of the Gross-Witten transition is quite
degenerate and appears to differ significantly from any
putative string theory dual of QCD \cite{ias-double-cut}.
On the other hand, there does
exist a more tangible string description of two dimensional QCD
due to Gross and Taylor \cite{grosstay15} which also has sigma
model representations \cite{moorhorav155}. The large $N$
transition in Wilson loops of continuum two dimensional QCD on
infinite space-time is of the same type as the Douglas-Kazakov
transition in the partition function of planar two dimensional QCD
defined on a sphere of area $A$ \cite{dougkaz157}.

There is another feature in favor
of two dimensional QCD as the universality
class and it has to do with instantons. When one takes the
nonabelian instanton solution that we showed makes a non-zero
contribution to the eigenvalue density at $-1$ for small loops in
four dimensions and looks only at a two dimensional slice
one gets two dimensional instantons that were
argued to be related to the Douglas-Kazakov transition
in~\cite{gross-mat}.

All this adds up to putting the Gross-Witten transition in
disfavor as a representative of the universality class of our
transitions, and leaves us with the preferred option of two
dimensional planar QCD. An explicit formula for the
eigenvalue distribution of Wilson loops can be found in~\cite{duol}
and it is possible~\cite{inst2d}
to connect this formula to (\ref{twodwil}).

\begin{figure}[ht]
\centering
\includegraphics[width=\textwidth]{images/durolefit.pdf}
\caption{ Fit of the
distributions to the Durhuus-Olesen distributions for four
different sizes of Wilson loops, namely,
$lt_c=0.740,0.660,0.560,0.503$. The associated areas ($k$ in the
Durhuus-Olesen notation) that describe the continuous curves are
given by $k=4.03,2.30,1.41,1.15$ respectively. }
\label{durolefit}
\end{figure}


To see that we are not completely off in our guess we show a plot
of several eigenvalue distributions we measure
versus the exact infinite $N$
continuum solution of Durhuus and Olesen \cite{duol}:

Figure~\ref{durolefit} shows shows how the critical size is
traversed and one sees directly the effect of the transition and
that the Durhuus-Olesen distribution seems to work pretty well on
both sides of the transition. An attempt to fit to the
Gross-Witten distribution would fail, since on the gap-less side
the Gross-Witten distribution would always cross the value
$\frac{1}{2}$ in Figure~\ref{durolefit}
at $\theta=\pm\frac{\pi}{2}$.


Another point in favor of the universality class is evident from
equation~\ref{twodwil}: The prefactor has oscillating signs for
large areas. These are nothing but the oscillating signs we saw in
the Fourier coefficients of Figure~\ref{fourier}. These
oscillating signs could be taken as an indication that an
effective string representation of loops traversed multiple times
requires some elementary two dimensional excitation of statistics
different from that of bosons~\cite{gross-ooguri}. The simple
effective string models one typically uses do not address loops
traversed multiple numbers of times.


Because of the importance of the large $N$ universality class to our
program, we now turn to a more general description of what appears to
be its fundamental representative family in terms of matrix models:

Imagine $n$ $SU(N)$ matrices, independently and identically
distributed according to a conjugation invariant measure
$d\mu(U_i)$. The measure is peaked at $U_i=1$ and is of the single
trace type. As a function of each eigenvalue it decreases
monotonically as one goes round the circle from $1$ to $-1$, where
it has a minimum. The continuum Wilson loop operators, $\hat W$,
one defines, would behave in the vicinity of the large $N$
transition point at which their eigenvalue distribution opens a
gap at $-1$, similarly to $W=\prod_i^n U_i$. The parameters of the measure
and $n$ can be adjusted to various critical behaviors in the
eigenvalue distribution of $W$ at infinite $N$, and likely produce
a variety of double-scaling limits corresponding to large $N$
fixed points of decreasing degrees of stability. The shape and
scale of ${\hat W}$ map analytically into these parameters in the
vicinity of the transition. Most plausibly, the most generic and
stable transition is the relevant one.

This random matrix model is exactly soluble, as we now sketch,
leaving details for the future:

One introduces $n$ pairs of $N$-component Grassmann variables,
$\bar\psi_i, \psi_i$ and proves by simple recursion that
\be
\int \prod_i^n ( d\bar\psi_i d\psi_i ) e^{z\sum_i^n
\bar\psi_i\psi_i - \sum_i^n \bar\psi_i U_i\psi_{i+1}}=\det(z^n + W )
\ee
Above, we use the convention $\psi_{n+1}\equiv - \psi_1$.

Now, one can introduce auxiliary
fields for the $SU(N)$ invariant bilinears
$\bar\psi_i \psi_i$ and $\bar\psi_{i}\psi_{i+1}$. Averaging over
the $U_i$ variables can be done independently and the remaining
integral can be dealt with at infinite $N$ by ordinary saddle
point methods.


Choosing an axial gauge on an infinite two dimensional lattice, it
is easy to see that both on the lattice and in the continuum
\cite{duol} the non-minimal Wilson loops of two dimensional planar
QCD are indeed described by the above universality class.
Similarly, the Douglas-Kazakov critical point is also in this
universality class, corresponding to the most stable critical
point.

Intuitively, the universality class can be visualized as a
nonabelian generalization of the following abelian situation:
Let there be given a smooth loop $C$ and let $S$ be a surface
of the topology of a disk bounded by $C$. On the surface $S$
define a two dimensional  $U(1)$ gauge field, with independently
and identically distributed fluxes through small patches of
surface that tile  $S$. The phase of the parallel transport
round $C$ is then given by the sum of all these fluxes. The
distribution of that sum will be generically governed by the
central limit theorem. It is now obvious that we have a nonabelian
generalization of this arrangement~\cite{janik}, of the type studied by
Voiculescu \cite{voicu}.
